\begin{figure*}[p]
\centering

\begin{promptbox}[Response Similarity Judge (Part 1/3)]
Rate how similar and complete each generated response is compared with the ground truth response.

You will be given past messages for [HUMAN], the current context, the ground truth response, and generated response(s) that you should evaluate.

Provided Information:
{context}

<|The Start of Ground Truth Response|>
{ground_truth}
<|The End of Ground Truth Response|>

{generations_text}

For each generated response, assign three scores in [0.0, 1.0]:

1. "semantic_similarity": how similar the generated response is to the ground truth response in core meaning and information.
Scoring guidelines:
1.0 = Identical meaning, all key information preserved
0.75-0.95 = Same core meaning with minor semantic differences
0.55-0.75 = Similar meaning but some information changes
0.35-0.55 = Related topics but different focus or emphasis
0.15-0.35 = Some semantic overlap but different main message
0.0-0.15 = Completely different meaning or topic

2. "information_completeness": how completely the generated response preserves and covers the information from the ground truth response.
Scoring guidelines:
1.0 = All information fully preserved
0.75-0.95 = Minor details missing but all key info present
0.55-0.75 = Some important information missing
0.35-0.55 = Significant information gaps
0.15-0.35 = Major information missing, only basics preserved
0.0-0.15 = Almost no information preserved

3. "score": a strict final overall alignment score based on semantic similarity and information completeness. The final score should not be high unless both semantic similarity and information completeness are high. Penalize unsupported extra content, irrelevant speculation, excessive verbosity relative to the ground truth, wrong perspective, and source copying.

Penalty guidance:
- Hard-penalize sycophantic or generic agreement responses, especially formulaic openings such as "you're absolutely right", "totally", "completely agree", "100%", "that is right", or "yeah".
- Hard-penalize generic compromise responses that soften, reframe, or dilute a specific ground-truth position into a broad, moderate, or generic takeaway, for example turning a concrete argument into "it's a good starting point, but not a silver bullet" or "the real issue is...".
- Topical relevance is not enough: if the response mainly adds plausible-sounding background, generic commentary, or reusable high-level framing without matching the ground truth's goals, values, communication style, beliefs, and emotions, assign a penalty.
\end{promptbox}

\vspace{-6pt}

\caption{Exact response similarity judge prompt used in the experiments, split across three parts for typesetting.}
\label{fig:sim_prompt}
\end{figure*}

\begin{figure*}[p]
\centering

\begin{promptbox}[Response Similarity Judge (Part 2/3)]
Examples for semantic similarity:

Reference: "I need to reschedule our meeting to Thursday at 3pm because I have a doctor's appointment on Tuesday."
Predicted: "I have a medical appointment on Tuesday, so can we move our meeting to Thursday afternoon at 3?"
Explanation: The predicted dialogue conveys the same core meaning as the reference with only very minor wording differences. All key information is preserved including the reason for rescheduling, the conflicting appointment, and the new proposed time. The slight paraphrasing does not change the semantic content.
<score>0.95</score>

Reference: "I'm excited about the new project because it involves machine learning and will help improve customer experience."
Predicted: "The new project is interesting and uses AI technology."
Explanation: While both dialogues discuss a new project involving AI/machine learning, the predicted version has a different emphasis and is missing important nuances. The reference conveys excitement and specifically mentions the goal of improving customer experience, whereas the predicted dialogue only expresses mild interest and uses the more general term "AI technology." The core topic is related but the focus and emotional tone differ significantly.
<score>0.50</score>

Reference: "I'm planning to visit Paris next summer to see the museums and try French cuisine."
Predicted: "My favorite programming language is Python because it's versatile and easy to learn."
Explanation: These two dialogues are about completely different topics with no semantic overlap whatsoever. The reference discusses travel plans to Paris, while the predicted dialogue is about programming language preferences. There is no connection in meaning or information between them.
<score>0.0</score>

Examples for information completeness:

Reference: "The patient should take 500mg of amoxicillin twice daily for 10 days. Avoid alcohol and dairy products within 2 hours of each dose."
Predicted: "Take 500mg amoxicillin twice daily for 10 days. Avoid alcohol and dairy around the time you take it."
Explanation: The predicted dialogue preserves all critical information from the reference including the exact dosage (500mg), frequency (twice daily), duration (10 days), and both restrictions (alcohol and dairy). The only minor difference is the less precise "around the time you take it" instead of "within 2 hours," but the core medical information is fully intact.
<score>0.95</score>

Reference: "Our company was founded in 1987 by Sarah Chen in Boston. We now have 500 employees across 12 offices and generated $50M in revenue last year."
Predicted: "The company was started by Sarah Chen and has grown to multiple offices with hundreds of employees."
Explanation: The predicted dialogue retains the founder's name but loses significant specific information. The founding year (1987) and location (Boston) are missing. The precise employee count (500) is reduced to vague "hundreds," the exact office count (12) becomes "multiple," and the revenue figure ($50M) is completely omitted. While the general narrative is preserved, most quantitative details are lost.
<score>0.50</score>

Reference: "The experiment showed a 23% increase in cell growth when exposed to 450nm blue light for 6 hours at 37°C, with p<0.001 significance."
Predicted: "The experiment showed that blue light increased cell growth."
Explanation: The predicted dialogue only preserves the most basic conclusion that blue light increased cell growth. All quantitative and methodological details are lost: the specific percentage (23%), wavelength (450nm), exposure duration (6 hours), temperature (37°C), and statistical significance (p<0.001). Only the fundamental relationship between blue light and cell growth remains.
<score>0.25</score>
\end{promptbox}
\end{figure*}

\begin{figure*}[p]
\centering

\begin{promptbox}[Response Similarity Judge (Part 3/3)]
Remember to judge each generated response against the ground truth response only. Do not compare generated responses to each other.

Perspective rule:
- Determine whether each generated response is written as [HUMAN].
- If it speaks as another user in the context, or treats another user's first-person statements as if they were [HUMAN]'s own, set "wrong_perspective" to true and set all three scores to 0.0.
- If the ground truth is [HUMAN]'s next question, follow-up, or short request, but the generated response instead behaves like an assistant reply by directly answering the earlier prompt, giving a polished explanatory answer, or providing a structured helpful breakdown, set "wrong_perspective" to true and set all three scores to 0.0.
- Otherwise set "wrong_perspective" to false.

Be strict and reserve scores above 0.8 for clearly outstanding matches.

Output format (JSON):
{
  "key_points": "<brief single-line analysis of the key information and meaning in the ground truth response>",
  "1": {"thought": "<single-line reasoning for the scores>", "semantic_similarity": <float in [0,1]>, "information_completeness": <float in [0,1]>, "score": <float in [0,1]>, "wrong_perspective": <true or false>, "source_copy": <true or false>},
  "2": ...
}

Format Notes:
- All text in "key_points" and "thought" fields MUST be on a single line with no line breaks or newlines.
- Use standard JSON string format with double quotes. For any quotes needed inside strings, use single quotes (').
- Ensure that all numbered items from 1 to {num_generations} are present.
- Ensure that every item includes "thought", "semantic_similarity", "information_completeness", "score", "wrong_perspective", and "source_copy".

Your output:
\end{promptbox}
\end{figure*}
